
\label{sec:coordenadas_esfericas}


\textbf{Cambio de coordenadas Lineales} \textcolor{red}{Agregar}

Para el cálculo de integrales en regiones con simetría esf\'ericas, se utiliza el sistema de coordenadas esf\'ericas. Un punto $(x, y, z)$ se relaciona con $(\rho, \theta ,  \phi)$ mediante:
\begin{equation}\label{def:coord_esfericas}
\mathbf{T}:W^*=[0,R]\times[0, 2\pi]\times[0,\pi]  \subseteq\Rn{3} \to W\: : \: \mathbf{T}(r, \theta,z ) =( \rho \sin \phi \cos \theta  ,  \rho \sin \phi \sin \theta ,  \rho \cos \phi ).
\end{equation} 

\begin{tarea} Probar que la transformaci\'on $T$ definida en (\ref{def:coord_esfericas}) est\'a bajo las las hi\'otesis del Teorema  ~\ref{teo:cambio_variable}
\end{tarea}


Para el c\'alculo del \emph{Determinante Jacobiano} de $T$  primero  calculamos su  matriz diferencial $D_\mathbf{T}$, quedando
 
\[
    \boldsymbol{D}_{\mathbf{F}}(\rho, \theta, \phi) 
    =
    \begin{pmatrix}
    \sin\phi \cos\theta & -\rho \sin\phi \sin\theta & \rho \cos\phi \cos\theta \\
    \sin\phi \sin\theta & \rho \sin\phi \cos\theta & \rho \cos\phi \sin\theta \\
    \cos\phi & 0 & -\rho \sin\phi
    \end{pmatrix}
    \]
 
 por \'ultimo,  en nuestro caso,  como $  0 \leq \phi \leq \pi$, tenemos
\begin{equation}   \label{ej:jacobiano_esferico}
    \left| J_{\mathbf{T}} \right| = \rho^2 \sin \phi
\end{equation}
  
\textbf{Visualización del sistema de coordenadas esféricas}

\begin{figure}[H]
    \centering
    % Llamamos al archivo con el código TikZ extraído
    
    \centering
    \tikzsetnextfilename{esfericas_cambio_figura}

    \begin{tikzpicture}

    %%% IZQUIERDA: caja W* en espacio (ρ, φ, θ)
    \begin{axis}[
        name=izq,
        view={110}{25},
        width=7cm,
        height=7cm,
        xmin=-0.2, xmax=3.5,
        ymin=-0.2, ymax=7.5,
        zmin=-0.2, zmax=3.8,
        xlabel={$\rho$},
        ylabel={$\theta$},
        zlabel={$\phi$},
        axis lines=center,
        ticks=none,
        every axis x label/.style={at={(ticklabel* cs:1.05)}, anchor=north east},
        every axis y label/.style={at={(ticklabel* cs:1.05)}, anchor=north west},
        every axis z label/.style={at={(ticklabel* cs:1.05)}, anchor=south},
    ]

    % Cara inferior (φ = 0)
    \addplot3[fill=magenta!20, opacity=0.85, draw=magenta!70, thick]
        coordinates {(0,0,0)(3,0,0)(3,6.28,0)(0,6.28,0)(0,0,0)};
    % Cara posterior (θ = 0)
    \addplot3[fill=magenta!15, opacity=0.7, draw=magenta!70, thick]
        coordinates {(0,0,0)(3,0,0)(3,0,3.14)(0,0,3.14)(0,0,0)};
    % Cara lateral (ρ = 0)
    \addplot3[fill=magenta!20, opacity=0.7, draw=magenta!70, thick]
        coordinates {(0,0,0)(0,6.28,0)(0,6.28,3.14)(0,0,3.14)(0,0,0)};
    % Cara frontal (ρ = R)
    \addplot3[fill=magenta!30, opacity=0.8, draw=magenta!70, thick]
        coordinates {(3,0,0)(3,6.28,0)(3,6.28,3.14)(3,0,3.14)(3,0,0)};
    % Cara lateral (θ = 2π)
    \addplot3[fill=magenta!25, opacity=0.75, draw=magenta!70, thick]
        coordinates {(0,6.28,0)(3,6.28,0)(3,6.28,3.14)(0,6.28,3.14)(0,6.28,0)};
    % Cara superior (φ = π)
    \addplot3[fill=magenta!35, opacity=0.9, draw=magenta!70, thick]
        coordinates {(0,0,3.14)(3,0,3.14)(3,6.28,3.14)(0,6.28,3.14)(0,0,3.14)};

    % Etiqueta W*
    \node[color=magenta!80!black, font=\Large]
        at (axis cs: 1.5, 3.14, 1.57) {$W^*$};

    % Cotas
    \node[below, color=black!60, font=\small]
        at (axis cs: 3, 0, 0) {$R$};
    \node[above right, color=black!60, font=\small]
        at (axis cs: 0, 6.28, 0) {$2\pi$};
    \node[left, color=black!60, font=\small]
        at (axis cs: 0, 0, 0) {$0$};
    \node[left, color=black!60, font=\small]
        at (axis cs: 0, 0, 3.14) {$\pi$};

    \end{axis}

    % Label inferior izquierdo
    \node[below=0.15cm of izq.south, font=\small]
        {Espacio $(\rho,\theta,\phi)$};

    %%% FLECHA T
    \draw[-{Stealth[length=9pt, width=6pt]}, thick]
        ($(izq.east)+(0.4,0)$)
        -- node[above, font=\large]{$T$}
        ++(1.5,0);

    %%% DERECHA: esfera W en espacio (x, y, z)
    \begin{axis}[
        name=der,
        at={($(izq.east)+(2.4cm,0)$)},
        anchor=west,
        view={110}{25},
        unit vector ratio=1 1 1,
        width=7cm,
        height=7cm,
        xmin=-3, xmax=3,
        ymin=-3, ymax=3,
        zmin=-3, zmax=3,
        xlabel={$x$},
        ylabel={$y$},
        zlabel={$z$},
        axis lines=center,
        ticks=none,
        every axis x label/.style={at={(ticklabel* cs:1.05)}, anchor=north east},
        every axis y label/.style={at={(ticklabel* cs:1.05)}, anchor=north west},
        every axis z label/.style={at={(ticklabel* cs:1.05)}, anchor=south},
    ]

    % Hemisferio inferior
    \addplot3[surf, draw=green!45,
              point meta=0, colormap={verdeinf}{color=(green!20) color=(green!20)},
              fill opacity=0.32, draw opacity=0.75,
              domain=0:360, y domain=90:180,
              samples=48, samples y=22, shader=interp]
        ({2*cos(x)*sin(y)}, {2*sin(x)*sin(y)}, {2*cos(y)});

    % Hemisferio superior
    \addplot3[surf, draw=green!55,
              point meta=0, colormap={verdesup}{color=(green!30) color=(green!30)},
              fill opacity=0.38, draw opacity=0.85,
              domain=0:360, y domain=0:90,
              samples=48, samples y=22, shader=interp]
        ({2*cos(x)*sin(y)}, {2*sin(x)*sin(y)}, {2*cos(y)});

    % Circulo ecuatorial (z = 0)
    \addplot3[draw=green!70, thick, domain=0:360, samples=60]
        ({2*cos(x)}, {2*sin(x)}, {0});

    % Etiqueta W
    \node[color=green!50!black, font=\Large]
        at (axis cs: 0.4, 0.4, 0.8) {$W$};

    \end{axis}

    % Label inferior derecho
    \node[below=0.15cm of der.south, font=\small]
        {Espacio $(x,y,z)$};

    \end{tikzpicture}

    \caption{Visualización del sistema de coordenadas esféricas. 
    El bloque $W^*$ a la izquierda representa la región en el espacio 
    $(\rho,\theta,\phi)$ con sus límites típicos, mientras que la esfera 
    $W$ a la derecha muestra cómo se transforma esa región en el 
    espacio $(x,y,z)$ mediante la transformación $T$.\href{https://www.geogebra.org/m/pzujyjqy}{Ver en GeoGebra }}
 
    \label{fig:Visualizacion esfericas}
\end{figure}

\vspace{1cm}
